% Canonical Paper IV section.
The first calibration stays inside the affine sector of AEG.  It shows exactly
how a commutative charge shadow can forget positional information while the
evaluated operators remain distinct.

\subsection{Radix histories}

Fix an integer radix $k\ge2$ and the digit alphabet
$D_k=\{0,1,\ldots,k-1\}$.  For $d\in D_k$, let
\[
  h_d(x)=kx+d.
\]
A chronological word $w=d_1d_2\cdots d_n$ acts by
\[
  \rho(w)=h_{d_n}\circ\cdots\circ h_{d_1}.
\]

\begin{proposition}[Horner evaluation formula]
\label{prop:horner-evaluation}
For every length-$n$ word,
\begin{equation}
  \rho(w)(x)
  =k^n x+\sum_{i=1}^{n}d_i k^{n-i}.
  \label{eq:horner-evaluation}
\end{equation}
\end{proposition}

\begin{proof}
Induction on $n$.  Appending $d_{n+1}$ sends the current value $y$ to
$ky+d_{n+1}$, multiplying all previous coefficients by $k$ and adding the new
least significant digit.
\end{proof}

\begin{theorem}[Radix-word injectivity]
\label{thm:radix-word-injectivity}
Over $\Z$, or over a characteristic-zero field containing the displayed
integers, distinct words in $D_k^n$ induce distinct affine operators.
\end{theorem}

\begin{proof}
All operators have slope $k^n$.  Equality would force equality of the constants
in Equation~\eqref{eq:horner-evaluation}.  These constants are the integers
$0,1,\ldots,k^n-1$ written uniquely in length-$n$ base-$k$ notation, allowing
leading zeros.
\end{proof}

The characteristic-zero hypothesis is active.  Over a finite field, distinct
integer constants may coincide after reduction; Section~\ref{sec:butterfly-ntt}
uses such modular identification deliberately.

\subsection{Condensation by ACS charge}

Realize $h_d$ as a scale by $k$ followed by an addition of $d$.  The ACS of
Paper~I~\cite{Yuan2026AEGFoundations} records only total additive and
logarithmic multiplicative charge.  Hence
\begin{equation}
  \chi(w)=\left(\sum_{i=1}^{n}d_i,\ n\log k\right).
  \label{eq:horner-charge}
\end{equation}

\begin{theorem}[Fixed-histogram operator count]
\label{thm:horner-histogram-count}
Let $m_r$ be the number of occurrences of digit $r$, with
$\sum_{r=0}^{k-1}m_r=n$.  All words with this histogram have the same ACS
charge, and they induce exactly
\begin{equation}
  \frac{n!}{\prod_{r=0}^{k-1}m_r!}
  \label{eq:multinomial-horner-count}
\end{equation}
distinct affine operators.
\end{theorem}

\begin{proof}
Equation~\eqref{eq:horner-charge} depends only on the histogram.  The number of
words with the histogram is the multinomial coefficient.  They give distinct
operators by Theorem~\ref{thm:radix-word-injectivity}.
\end{proof}

\begin{corollary}[Binary fixed-weight fiber]
\label{cor:binary-horner-count}
For $k=2$, the length-$n$ words with exactly $m$ ones have one ACS charge and
$\binom nm$ distinct operator images.  Under the uniform distribution on that
set, conditioning only on the charge loses
\[
  \log_2\binom nm
\]
bits of word information.
\end{corollary}

For $n=4$ and $m=2$, the six constants are
\[
  3,5,6,9,10,12,
\]
corresponding to the six placements of two ones.  The charge remembers their
number but not their positions.

\subsection{Which state is sufficient?}

For all future affine continuations, the current affine operator is sufficient:
composition depends only on its slope and translation.  The ACS charge is not
sufficient, because two words with one charge can already have different
translations before any future is applied.  If the observation is deliberately
coarsened to ACS charge and all future updates act only on charge, then the
charge becomes an exact residual for that smaller experiment.

This illustrates the hierarchy
\[
  \text{word}
  \longrightarrow
  \text{affine operator}
  \longrightarrow
  \text{ACS charge}.
\]
The size in Theorem~\ref{thm:horner-histogram-count} is the operator image of a
charge fiber.  It is not a fiber of operator evaluation.

\subsection{Streaming cost and encoding cost}

The recurrence
\[
  a_i=ka_{i-1}+d_i
\]
evaluates a word in $n$ affine steps while storing one current field element.
In a unit-cost field-operation model this is $O(n)$ work and $O(1)$ field
elements of workspace.  Over the integers, however, the bit length of $a_i$
is $\Theta(i\log k)$ in the worst case.  Thus ``one accumulator'' does not mean
constant bit space in a bit-complexity model.

The example therefore supplies an exact condensation count but no generic
search or hardness theorem.  Positional information can be large while direct
streaming evaluation remains simple.
